% s-219-p219-m002; original PDF 219, 220
\Aphorism{s-219-p219-m002}{312}
Das kleine Getriebe, das Spiel der feinsten Räder und Federn ist unter Umständen Zweck der psychologischen Dichtung. Bisweilen aber auch das Gegentheil, und dann gebietet die künstlerische Öconomie, die Aufmerksamkeit von den grossen Motoren, die ausgestellt werden, nicht in ihr gar zu minutiöses Detail hinein abzulenken, mit anderen Worten, \Sperr{keine} überflüssige Psychologie zu treiben. Die richtige Function der Organe und Mechanismen niederer Ordnung wird dann vorausgesetzt; zur Einsicht geöffnet und beleuchtet sind nur die reicheren, zusammengesetzteren Gebilde. Freilich kennt die Natur kein Höher oder Niedriger, kein Einfacher und Complicirter; die Kunst bringt also eine fremde, auf den Zuschauer berechnete Rangordnung in die Dinge, eine Exposition wie für den Schauladen, die mit vieler Kunst vergessen gemacht werden muss. Die Bequemlichkeit des Sehens darf selbst nicht ahnen, wie sehr ihr geschmeichelt wird; man erhält die Natur präsentirt, muss aber wähnen, sie zu überraschen.
\par

